\begin{abstract}
Verifiable code generation asks a language model for code that a verifier can prove correct against a formal specification. The VERINA benchmark showed that this is hard in Lean, but it evaluated models in a single shot, in one language and without a test-driven baseline. We study what happens when the same tasks are given to an agent that receives feedback. We re-expressed \NTASKS{} VERINA-BASIC tasks in Dafny and in Python with iContract contracts checked by CrossHair, and we checked that the two versions of each task are equivalent using shared tests, specification probes and differential testing. We then built SpecLoop, an agent whose feedback sources can be switched on and off, with a router that uses test results to separate wrong programs from missing proofs and a guard that rejects any change to the specification. With GPT-4o-mini and three seeds per configuration, VERINA's protocol passes \R{etwo.dafny.pass}\% of tasks in Dafny and \R{etwo.python.pass}\% in Python, against the \R{etwo.verina.pass}\% reported in Lean. On Dafny, testing-only feedback raises correctness from \R{dafny.single_shot.correct}\% to \R{dafny.testing_only.correct}\% but stops with an unverified program in \R{dafny.testing_only.falsesuccess}\% of runs, verifier feedback reaches \R{dafny.verifier_only.correct}\%, and SpecLoop reaches \R{dafny.specloop_full.correct}\%. Python is near the ceiling for every configuration. The results suggest that the choice of specification language decides most of the outcome, that a verifier is worth adding to a test-driven agent mainly because it tells the agent when it is not done, and that tests make verifier feedback easier to act on. All code, data, trajectories and model responses are released so that every number can be reproduced without API calls.
\end{abstract}

\begin{IEEEkeywords}
LLM agents, verifiable code generation, formal verification, Dafny, design by contract, program repair
\end{IEEEkeywords}
